% Generated draft: requires mathematical and source review.
\documentclass[11pt]{article}
\usepackage[margin=25mm]{geometry}
\usepackage{fontspec}
\setmainfont{Latin Modern Roman}
\usepackage{amsmath,amssymb,mathtools,unicode-math}
\setmathfont{Latin Modern Math}
\usepackage{enumitem,needspace,xurl,hyperref}
\setlength{\parindent}{0pt}
\setlength{\parskip}{5pt}
\setlength{\emergencystretch}{3em}
\newcommand{\N}{\mathbb{N}}
\newcommand{\Z}{\mathbb{Z}}
\newcommand{\Q}{\mathbb{Q}}
\newcommand{\R}{\mathbb{R}}
\newcommand{\C}{\mathbb{C}}
\newcommand{\F}{\mathbb{F}}
\newcommand{\A}{\mathbb{A}}
\newcommand{\PP}{\mathbb{P}}
\newcommand{\E}{\mathbb{E}}
\newcommand{\eps}{\varepsilon}
\newcommand{\dd}{\,\mathrm{d}}
\begin{document}
\section*{UnsolvedMath ID: 30005342}
\textbf{Source problem code:} catalogue entry
\subsection{Definitions and mathematical statement}
For a finite simple undirected graph \(G=(V,E)\), a proper vertex coloring assigns colors to vertices so that adjacent vertices receive distinct colors. The chromatic number \(\chi(G)\) is the minimum number of colors in such a coloring. The maximum degree is \(\Delta(G)=\max_{v\in V}\deg_G(v)\), defined here for nonempty \(V\), and the clique number \(\omega(G)\) is the largest cardinality of a set of pairwise adjacent vertices. For a real number \(x\), \(\lceil x\rceil\) denotes the least integer greater than or equal to \(x\).

The supplied clean statement asks whether every graph \(G\) satisfies
\[
\chi(G)\leq \left\lceil\frac{\Delta(G)+1+\omega(G)}{2}\right\rceil.
\]
Under the explicit, provisional interpretation that the graphs are finite, simple, undirected, and nonempty, the candidate statement is
\[
\forall\,G\text{ finite, simple, undirected, with }V(G)\ne\varnothing,\qquad
\chi(G)\leq \left\lceil\frac{\Delta(G)+1+\omega(G)}{2}\right\rceil.
\]
The supplied record does not explicitly specify this graph domain or a convention for the empty graph. Consequently, this precise candidate is conditional on confirmation of those domain choices; it does not settle the meaning of the unrestricted phrase ``every graph.''

The record titles the problem Reed's Chromatic-Number Conjecture and labels it open. Its clean statement is marked \texttt{corrected\_verified}, with a supplied claim of verification against the official report. However, the retained original passage does not state the displayed inequality, and the supplied dated literature assessment says that no identifiable mathematical assertion was extracted. These provenance claims require reconciliation; neither verification nor open status is independently established here.
\subsection{Short English statement}
The supplied question asks whether every graph can be properly vertex colored using at most the ceiling of half the sum of its maximum degree, its clique number, and one. A precise candidate is given for finite simple undirected nonempty graphs, but the intended graph domain and the connection between the clean statement and the retained report passage require review.
\subsection{Sources}
\begin{itemize}
\item Local UnsolvedMath registry, permanent ID \(30005342\), source problem code \texttt{catalogue entry}; supplied title: ``Reed's Chromatic-Number Conjecture.'' The registry supplies the displayed inequality, the status \texttt{open}, and the statement status \texttt{corrected\_verified}. No proposer is specified; the supplied proposed year is 2023.
\item Supplied report citation: \emph{Combinatorics} (2023), Oberwolfach Reports, pp.~5--89, DOI \texttt{10.4171/owr/2023/1}. Supplied source and literature link, labeled ``Oberwolfach Reports 2023/1'': \url{https://doi.org/10.4171/owr/2023/1}. The record describes Oberwolfach Reports as documenting mathematical workshops and their open-problem discussions. No precise location within the report is supplied for the displayed inequality.
\item Supplied literature assessment dated \texttt{2026-08-22}: the item is classified open because the assessment describes the exact record as containing only the label ``Conjecture'' and no identifiable mathematical assertion. This is supplied dated curation, not an independently verified literature assessment. The record itself requests independent verification of time-sensitive status and cited claims.
\item The retained original passage mentions Mozhan partitions and move sequences, attributes a comparison argument to Cranston and Rabern in reference [7], refers to an unspecified Theorem 1, and says that the case \(t=1\) was treated in reference [17] with a clique-size lower bound of \(\Delta(G)-17\). The bibliography entries, theorem statement, parameter definition, and applicable hypotheses are not supplied.
\end{itemize}
\end{document}
